% Part: proof-theory
% Chapter: natural-deduction
% Section: introduction

\documentclass[../../../include/open-logic-section]{subfiles}

\begin{document}

\olfileid{pt}{nat}{int}

\olsection{مقدمه}

استنتاج طبیعی خانواده‌ای از دستگاه‌های !!{proof} است که در آن هر رابط
منطقی یا سور، یک جفت استنتاج دارد: یکی عملگر را «معرفی» می‌کند
(قاعدهٔ !!a{introduction}) و دیگری عملگر را «حذف» می‌کند
(قاعدهٔ !!a{elimination}). برای نمونه، نتیجهٔ \Intro{\lif} یک
!!{formula} به صورت~$!A \lif !B$ است، و مقدمهٔ \Elim{\lif} یک
!!{formula} به صورت~$!A \lif !B$.

نخستین دو صورت استنتاج طبیعی را در دههٔ ۱۹۳۰ گرهارت گنتسن و
استانیسواف یاشکوفسکی معرفی کردند. نظریه‌پردازان اثبات بیشتر دستگاه
گنتسن را بررسی می‌کنند، و دستگاه‌های رایج‌تر در آموزش از کار یاشکوفسکی
پدید آمده‌اند. در هر دو دستگاه می‌توان از \emph{فرض‌ها} آغاز کرد؛ یعنی
!!{formula}هایی که لازم نیست !!{prove} شوند، اما می‌توان آن‌ها را در
!!{proving} دیگر !!{formula}ها به کار برد. ممکن است کل یک !!{proof} به
چنین فرض‌هایی وابسته باشد، و چون آن فرض‌ها اثبات نشده‌اند، !!{proof}
نتیجهٔ خود، یعنی \emph{!!{formula} پایانی}، را به‌تنهایی ثابت نمی‌کند؛
فقط نشان می‌دهد نتیجه از فرض‌ها به دست می‌آید.

برخی قواعد استنتاج چنان‌اند که نتیجهٔ آن‌ها دیگر به بعضی فرض‌ها وابسته
نیست: آن فرض‌ها را «مرخص» می‌کنند. برای نمونه، قاعدهٔ \Intro{\lif} یک
!!a{proof} برای نتیجهٔ~$!B$ از فرض~$!A$ می‌گیرد و نتیجهٔ
$!A \lif !B$ را می‌دهد. !!{proof}ی که به~$!A \lif !B$ پایان می‌یابد
دیگر به~$!A$ وابسته نیست؛ فرض~$!A$ مرخص شده است. در دستگاه‌های گنتسن،
!!{proof}ها درخت‌هایی از !!{formula}ها هستند، فرض‌ها بالاترین
!!{formula}ها هستند، و قواعدی مانند~\Intro{\lif} برچسب‌هایی دارند که
نشان می‌دهند کدام فرض‌ها مرخص می‌شوند. در دستگاه‌های یاشکوفسکی،
بخش‌هایی از !!{proof} که به یک فرض وابسته‌اند به نحوی جدا می‌شوند؛
مثلاً پشت یک خط عمودی تورفته‌اند یا در کادری قرار می‌گیرند. فرض~$!A$
که مرخص می‌شود و تالی~$!B$ شرطی، نخستین و آخرین سطر در کادرند، و
نتیجهٔ~$!A \lif !B$ از \Intro{\lif} بیرون کادر قرار می‌گیرد.

این دستگاه‌ها «طبیعی»اند، زیرا قواعد استنتاج بازتاب شیوه‌ای هستند که
ما، یا دست‌کم ریاضی‌دانان، به‌طور طبیعی استدلال می‌کنیم. برای اثبات یک
نتیجهٔ فرضی، یعنی یک شرطی، مقدم را فرض می‌کنیم و از آن برای اثبات تالی
بهره می‌بریم؛ این همان~\Intro{\lif} است. وقتی یک شرطی
$!A \lif !B$ و مقدم آن~$!A$ را می‌دانیم،~$!B$ را نتیجه می‌گیریم؛ این
همان~\Elim{\lif} است. برای نشان‌دادن اینکه با فرض یک فصل
$!A \lor !B$ نتیجه‌ای برقرار است، از برهان برحسب حالات استفاده
می‌کنیم؛ این~\Elim{\lor} است. برای اثبات ادعای کلی
$\lforall[x][!A(x)]$ نشان می‌دهیم $!A(c)$ برای یک~$c$ دلخواه برقرار
است؛ این~\Intro{\lforall} است، و به همین ترتیب.

برای نمونه، این !!{proof} ساده‌ای از~$!A \lif (!A \lor !B)$ در استنتاج
طبیعی به سبک گنتسن است:
\[
\AxiomC{$\Discharge{!A}{x}$}
\RightLabel{\Intro{\lor}}
\UnaryInfC{$!A \lor !B$}
\DischargeRule{\Intro{\lif}}{x}
\UnaryInfC{$!A \lif (!A \lor !B)$}
\DisplayProof
\]
این برهان با فرض~$!A$ آغاز می‌شود، با \Intro{\lor} نتیجهٔ
$!A \lor !B$ را می‌گیرد و سپس با \Intro{\lif} به
$!A \lif (!A \lor !B)$ می‌رسد. استنتاج \Intro{\lif} فرض~$!A$ را
مرخص می‌کند؛ برچسب~$x$ این امر را نشان می‌دهد.

نظریه‌پردازان اثبات دستگاه‌های اصلی گنتسن را که با درخت‌هایی از
!!{formula}ها کار می‌کنند ترجیح می‌دهند، هرچند یک صورت دیگر نیز در
نظریهٔ اثبات مهم است. !!^a{proof} برای~$!A$ از فرض‌های~$\Gamma$ نشان
می‌دهد $!A$ از~$\Gamma$ به دست می‌آید، یعنی
$\Gamma \Entails !A$. قواعد استنتاج طبیعی را نیز می‌توان به‌طور طبیعی
حاوی اطلاعاتی دربارهٔ چنین پیامدهایی خواند؛ مثلاً \Intro{\lif} می‌گوید
اگر $\Gamma + !A \Entails !B$، آنگاه
$\Gamma \Entails !A \lif !B$. پس طبیعی است قواعد را مستقیماً چنان
صورت‌بندی کنیم که هم بر فرض‌ها و هم بر نتیجه‌ای عمل کنند که در مقدمات
و نتیجهٔ قاعدهٔ استنتاج از آن فرض‌ها به دست می‌آید؛ یعنی بر دنباله‌های
$\Gamma \Sequent !A$ عمل کنند. برهان سادهٔ بالا در چنین دستگاهی این
شکل را دارد:
\[
\Axiom$x:!A \fCenter !A$
\RightLabel{\Intro{\lor}}
\UnaryInf$x:!A \fCenter !A \lor !B$
\DischargeRule{\Intro{\lif}}{x}
\UnaryInf$\fCenter !A \lif (!A \lor !B)$
\DisplayProof
\]
همان‌طور که می‌بینید، قواعد همان‌اند: بر !!{formula} سمت راست
~$\Sequent$ عمل می‌کنند، و !!{formula}های سمت چپ فقط فرض‌هایی را فهرست
می‌کنند که هر گام به آن‌ها وابسته است.

مجموعهٔ کامل جفت‌قواعد !!{introduction}/!!{elimination} به‌تنهایی برای
منطق کلاسیک تمام نیست. باید دو قاعدهٔ دیگر دربارهٔ~$\lfalse$ بیفزاییم.
نخست «قاعدهٔ شهودگرایانهٔ محال»~\FalseInt، یا «انفجار»، است که می‌گوید
از $\lfalse$ می‌توان هر چیزی را نتیجه گرفت. دومی اصل کلاسیک برهان
خلف~\FalseCl{}، یا «قاعدهٔ کلاسیک محال»، است که می‌گوید اگر بتوانیم از
~$\lnot !A$، $\lfalse$ را !!{prove} کنیم،~$!A$ را !!{prove} کرده‌ایم.
اگر~\FalseCl را کنار بگذاریم، دستگاهی مناسب منطق شهودگرایانه داریم؛ و
اگر هر دو را کنار بگذاریم، آنچه «منطق مینیمال» نامیده می‌شود به دست
می‌آید.

استنتاج طبیعی به سبک گنتسن برای منطق کلاسیک و شهودگرایانه را بررسی
خواهیم کرد. این‌ها دستگاه‌های \Log{N1c} و \Log{N1i} هستند. دستگاه‌های
\Log{N2c} و \Log{N2i} دستگاه‌های متناظرند که به جای !!{formula}های
تکین~$!A$ با دنباله‌های $\Gamma \Sequent !A$ کار می‌کنند.

\end{document}
